\documentclass[11pt,a4paper]{ctexart}
\usepackage[a4paper,margin=22mm,headheight=15pt]{geometry}
\usepackage{amsmath,amssymb,booktabs,longtable,tabularx,array}
\usepackage{xcolor,hyperref,fancyhdr,enumitem,float,graphicx}
\usepackage{setspace}
\setstretch{1.20}
\setlist{nosep,leftmargin=2em}
\pagestyle{fancy}
\fancyhf{}
\lhead{A题药材烘干}
\rhead{代码说明与结果可行性分析}
\cfoot{\thepage}
\hypersetup{colorlinks=true,linkcolor=black,urlcolor=blue}
\newcommand{\QThreeDryHours}{57.5303}
\newcommand{\QFourDryHours}{52.6776}
\newcommand{\AppendixFourFixedHours}{129.9873}
\newcommand{\NoCoordinateCaseHours}{51.1141}
\newcommand{\ShrinkReduction}{59.47\%}
\newcommand{\AppendixIncrease}{125.95\%}
\newcommand{\OmitCoordinateDifference}{-2.97\%}
\newcommand{\NetReduction}{8.44\%}
\newcommand{\StableStartHours}{2.050}
\newcommand{\PlateauTemperature}{50.0049}
\newcommand{\PlateauMoisture}{0.049994}
\newcommand{\QFourDryRadius}{1.2000}
\newcommand{\ConservationError}{1.61e-15}
\newcommand{\EquilibriumError}{2.06e-11}
\newcommand{\QThreeGridChange}{0.092\%}
\newcommand{\QFourGridChange}{0.045\%}
\newcommand{\QFourTimeChange}{0.00148\%}
\newcommand{\QOneMoistureProfileError}{8.69e-05}
\newcommand{\QTwoMoistureProfileError}{2.92e-06}


\title{\bfseries 药材烘干数值求解代码说明书\\[0.4em]\large 运行、输出与结果可行性分析}
\author{}
\date{2026年9月}

\begin{document}
\maketitle

\section{交付内容}
支撑文件采用“输入数据---数值求解---独立验证---结果导出”的可复现结构：
\begin{itemize}
  \item \texttt{src/drying\_solver.py}：附件读取、环境和半径插值、物性函数、有限体积组装、BDF1/BDF2及Picard迭代；
  \item \texttt{run\_all.py}：顺序运行四问、终止事件、机制对照和收敛计算，生成CSV与\texttt{summary.json}；
  \item \texttt{validate\_results.py}：平衡态、封闭系统守恒、关键表格和网格时间收敛验证；
  \item \texttt{tools/build\_workbooks.mjs}：将中间CSV写入题目规定的四个Excel文件；
  \item \texttt{tools/generate\_report\_assets.py}、\texttt{tools/make\_figures.py}：自动生成论文表格、宏和图片。
\end{itemize}

\section{运行方法}
Python依赖只有NumPy、Pandas和Matplotlib。Excel导出使用支撑环境中的Artifact Tool。在项目根目录执行：
\begin{verbatim}
python -m pip install -r requirements.txt
python run_all.py
python validate_results.py
python tools/generate_report_assets.py
python tools/make_figures.py
node tools/build_workbooks.mjs
\end{verbatim}
默认正式参数为：问题1--3取641个径向节点，问题4取321个无量纲径向节点；问题1、2前3 h使用1 s步长，长期阶段使用30 s步长。运行 \texttt{python run\_all.py --help} 可查看可调参数。

\section{程序模型}
固定半径问题求解
\begin{align}
\rho c_pT_t&=\frac1r(rkT_r)_r,\\
C_t&=\frac1r(rDC_r)_r,
\end{align}
中心为零通量，表面为Robin通量。问题4用 $\xi=r/R(t)$ 映射到 $[0,1]$，主程序保留
\begin{align}
T_t&=\xi\frac{\dot R}{R}T_\xi+\frac{1}{\rho c_pR^2\xi}(\xi kT_\xi)_\xi,\\
C_t&=\xi\frac{\dot R}{R}C_\xi+\frac{1}{R^2\xi}(\xi DC_\xi)_\xi.
\end{align}
参考代码原先把省略 $\xi\dot R/R$ 项的结果作为问题4主值；依据分享对话中的完整坐标变换推导，本交付已将其修正为主模型保留几何输运项，省略项仅作为敏感性算例。

\section{数值实现}
\subsection{空间离散}
问题1--3首先在物理径向坐标上取
\[
r_i=i\Delta r,\qquad r_{i+1/2}=\frac{r_i+r_{i+1}}2;
\]
问题4则在固定相对坐标上取
\[
\xi_i=i\Delta\xi,\qquad \xi_{i+1/2}=\frac{\xi_i+\xi_{i+1}}2,
\qquad r_i(t)=R(t)\xi_i.
\]
也就是说，代码划分的是 $r$ 或 $\xi$，不是直接划分体积。对相邻半节点之间的径向区间作守恒积分后，才自然得到
\[
\omega_i^r=\frac12(r_{i+1/2}^2-r_{i-1/2}^2),\qquad
\omega_i^\xi=\frac12(\xi_{i+1/2}^2-\xi_{i-1/2}^2)
\]
这两个几何权重。内部扩散项统一写成左右界面通量之差，界面 $k,D$ 采用调和平均；相邻区间共享同一界面通量，求和时内部通量逐项抵消。圆心区间从0到第一个半节点，其左端通量因半径为0自然消失，避免计算 $1/r$；最外区间直接代入 $h(T_a-T_s)$ 或 $h_m(C_a-C_s)$。问题4的扩散项再乘 $R^{-2}(t)$，几何输运项在内部节点使用中心差分、表面使用后向差分。离散后每个扩散场形成三对角系统，使用Thomas算法求解。

\subsection{时间推进和迭代}
首步用后向欧拉，之后用常步长BDF2：
\[
\frac{3\phi^{n+1}-4\phi^n+\phi^{n-1}}{2\Delta t}=L(\phi^{n+1})+b^{n+1}.
\]
每个时间层用Picard迭代更新物性，收敛阈值 $2\times10^{-10}$，正式算例单步最多8次迭代。问题3、4监测 $\max_iC_i-0.15$ 的下降过零，并在线性插值后保存精确终止时刻剖面。

\section{结果文件}
\begin{table}[H]
\centering
\caption{题目要求的Excel输出}
\begin{tabularx}{\textwidth}{lclX}
\toprule
文件 & 工作表 & 时间分辨率 & 空间列\\
\midrule
result1.xlsx & 温度、水分浓度 & 1 s，至1800 s & 0--2 cm，每0.1 cm\\
result2.xlsx & 温度、水分浓度 & 1 s，至10800 s & 0--2 cm，每0.1 cm\\
result3.xlsx & Sheet1 & 60 s，至结束前最后整分钟 & 0--2 cm，每0.1 cm\\
result4.xlsx & Sheet1 & 60 s，至结束前最后整分钟 & 0--2 cm固定位置及动态表面\\
\bottomrule
\end{tabularx}
\end{table}
问题4中若固定物理位置 $r_j>R(t)$，该位置已经在药材外，相应单元格为空；最后一列“药材表面”始终取 $\xi=1$。所有数值以数值单元格写入，显示四位小数。

\section{主要结果}
问题3的全场达标时间为
\[
\boxed{t_3=\QThreeDryHours\ \mathrm h},
\]
问题4完整移动边界模型为
\[
\boxed{t_4=\QFourDryHours\ \mathrm h},
\]
对应半径为\QFourDryRadius\ cm。作为机制对照，附录4固定半径时为\AppendixFourFixedHours\ h；使用附件2收缩但省略几何输运项时为\NoCoordinateCaseHours\ h。后者相对完整问题4低约2.97\%，说明该项虽不改变“约2.2 d”的量级，却影响四位小数主结果。

\section{可行性验证}
\begin{table}[H]
\centering
\caption{数值验证指标}
\begin{tabular}{lr}
\toprule
检验 & 结果\\
\midrule
封闭系统水分总量相对误差 & \ConservationError\\
均匀平衡态最大绝对误差 & \EquilibriumError\\
问题3：641到1281节点时长变化 & \QThreeGridChange\\
问题4：321到641节点时长变化 & \QFourGridChange\\
问题4：30 s到15 s时间步变化 & \QFourTimeChange\\
问题1：321与641节点含水率表值最大差 & \QOneMoistureProfileError\\
问题2：321与641节点含水率表值最大差 & \QTwoMoistureProfileError\\
\bottomrule
\end{tabular}
\end{table}
此外，问题1温度在28.0000--36.7856 $^\circ$C内，含水率在1.5102--2.5500 kg/kg内；问题3、4终止时最大含水率均位于中心；问题4结束早于附件2的72 h数据上限。这些检查支持边界符号、守恒组装、终止条件和半径插值的正确性。

\section{已知限制}
模型忽略轴向传递、蒸发潜热和长度收缩；问题2--4沿用题目唯一给出的 $h,h_m$。附件2只给出 $R(t)$，没有收缩本构，因此问题4是外生几何驱动模型。若获得实验终点或全过程测量，应优先校准表面对流传质、平衡含水率与 $R(C)$ 关系，并用守恒型ALE格式复核移动网格项。

\end{document}
